\subsection{Decay of high-order coupling}
\label{sec:ch463}
% TOC tag: [HOME][USE SS3.5.3]

Mode density on \(S^{2}\) in Subsection~\ref{sec:ch353} counts how many angular channels may be labeled at cutoff \(\ell_{\max}\)
\cite{jackson1999,hansen1988}. Subsection~\ref{sec:ch463} converts that count into a realizability law on
\(\boldsymbol{\mathcal{K}}_{\mathrm{geo}}\) from Eq.~\eqref{eq:ch4.geometry-coupling-operator}: listable channels are not exportable
channels \cite{stratton1941,harrington2001,devaney2004}. The philosophical insight is that DOS measures coordinate cardinality; coupling tails
measure watts. Truncation is inevitable because tail mass decays exponentially while labels grow quadratically in \(\ell_{\max}\).

\paragraph{Assumptions.}
Fix compact \(\Omega_{s}\) with \(\operatorname{diam}=D_{s}\), flux gauge, and coupling operator \(\boldsymbol{\mathcal{K}}_{\mathrm{geo}}\)
\cite{stratton1941,hansen1988}. Import Eq.~\eqref{eq:ch4.coupling-ell-decay} from Section~\ref{sec:ch45} by reference
\cite{harrington2001,devaney2004}. Narrowband linear exterior \cite{jackson1999}.

\paragraph{Coupling tail envelope.}
For \(\ell>k_{0}D_{s}\),
\begin{equation}
\label{eq:ch4.coupling-tail-envelope}
\sup_{\Jimp\in\mathcal{J}_{\mathrm{real}}}
\frac{|\mathcal{K}_{\ell m}^{(\sigma)}[\Jimp]|}{\|\boldsymbol{\mathcal{K}}_{\mathrm{geo}}[\Jimp]\|}
\le
\mathcal{C}_{\mathrm{tail}}(\ell,k_{0}D_{s})\,
\ell^{-1}\exp(-\ell/k_{0}D_{s}),
\end{equation}
extending Eq.~\eqref{eq:ch4.coupling-ell-decay} with sector constants \(\mathcal{C}_{\mathrm{tail}}\)
\cite{hansen1988,stratton1941,jackson1999}. Mathematically, Eq.~\eqref{eq:ch4.coupling-tail-envelope} is a uniform operator bound on
high-\(\ell\) blocks. Physically, it states that sources on \(\Omega_{s}\) cannot strongly excite angular orders far beyond electrical size
\cite{harrington2001,devaney2004}.

\begin{theorem}[Coupling tail implies spectral compression]
\label{thm:ch4.coupling-tail-compression}
If Eq.~\eqref{eq:ch4.coupling-tail-envelope} holds, then
\(r_{\mathrm{rad}}(\varepsilon)\le C(k_{0}D_{s},\varepsilon)\) with \(C\) growing slower than raw channel count \((2\ell_{\max}+1)^{2}\)
from Subsection~\ref{sec:ch353} \cite{hansen1988,stratton1941,harrington2001}. Listing size overstates synthesis dimension when
\(\ell_{\max}\gg k_{0}D_{s}\) \cite{devaney2004}.
\end{theorem}
\begin{proof}
Summing Eq.~\eqref{eq:ch4.coupling-tail-envelope} over \((m,\sigma)\) yields exponentially small tail mass for \(\ell\gg k_{0}D_{s}\)
\cite{hansen1988}. Comparing summed mass to \(\varepsilon\sigma_{1}\) in Definition~\ref{def:ch4.radiative-rank} bounds
\(r_{\mathrm{rad}}(\varepsilon)\) independently of \(\ell_{\max}\) \cite{harrington2001,stratton1941}.
\end{proof}

\paragraph{Cumulative coupling tail mass.}
Define
\begin{equation}
\label{eq:ch4.coupling-tail-mass}
\mathcal{M}_{\mathrm{cpl}}(\ell_{\mathrm{cut}})=\sum_{\ell>\ell_{\mathrm{cut}}}(2\ell+1)\max_{m,\sigma}|\mathcal{K}_{\ell m}^{(\sigma)}|^{2}.
\end{equation}
\begin{proposition}[Coupling tail versus DOS]
\label{prop:ch4.coupling-tail-dos}
\(\mathcal{M}_{\mathrm{cpl}}(\ell_{\mathrm{cut}})\) decays exponentially for \(\ell_{\mathrm{cut}}\gg k_{0}D_{s}\) while DOS grows as
\(\mathcal{O}(\ell_{\max}^{2})\) \cite{hansen1988,stratton1941,harrington2001}.
\end{proposition}

\paragraph{Tail mass integral bound.}
Integrating the envelope gives \(\mathcal{M}_{\mathrm{cpl}}(\ell_{\mathrm{cut}})\lesssim\mathcal{C}_{\mathrm{tail}}\exp(-\ell_{\mathrm{cut}}/k_{0}D_{s})\)
\cite{hansen1988,stratton1941}. At \(\ell_{\mathrm{cut}}=\ell_{\mathrm{supp}}(\varepsilon)\),
\begin{equation}
\label{eq:ch4.tail-mass-bound}
\mathcal{M}_{\mathrm{cpl}}(\ell_{\mathrm{supp}})\le\varepsilon^{2}\sigma_{1}^{2},\qquad
\sum_{\ell\le\ell_{\mathrm{supp}}}(2\ell+1)\ge d_{\mathrm{eff}}(\varepsilon).
\end{equation}
Equation~\eqref{eq:ch4.tail-mass-bound} links angular cutoff to export threshold \cite{harrington2001,devaney2004}.

\begin{figure}[t]
\centering
\ChOneFigBlock{%
\begin{tikzpicture}[ch1fig, node distance=7mm and 9mm]
  \node[ch1box] (dos) {DOS $\sim\ell_{\max}^{2}$};
  \node[ch1box, right=11mm of dos] (tail) {Eq.~\eqref{eq:ch4.coupling-tail-envelope}};
  \node[ch1box, right=11mm of tail] (rr) {$r_{\mathrm{rad}}$};
  \draw[ch1flow] (dos) -- (tail) -- (rr);
\end{tikzpicture}%
}
\caption{Subsection~\ref{sec:ch463}: channel DOS grows with listing; coupling tails compress export rank.}
\label{fig:ch4.coupling-tail-compression}
\end{figure}

Figure~\ref{fig:ch4.coupling-tail-compression} is the conceptual compression map: designers must not equate channel labels with independent
beams \cite{hansen1988,harrington2001}.

\paragraph{Worked illustration (high channel count, low rank).}
At \(\ell_{\max}=20\), \(k_{0}D_{s}=1.5\), Subsection~\ref{sec:ch353} lists hundreds of labels but \(r_{\mathrm{rad}}(10^{-2})\approx10\)
by Theorem~\ref{thm:ch4.coupling-tail-compression} \cite{hansen1988,harrington2001,devaney2004}.

\paragraph{Worked illustration (tail mass fraction).}
At \(\ell_{\max}=18\), \(k_{0}D_{s}=1.2\), only \(\ell\le5\) carry \(99\%\) of \(\mathcal{M}_{\mathrm{cpl}}\) \cite{hansen1988,devaney2004}.

\paragraph{Worked numerics (cutoff sweep).}
At \(k_{0}D_{s}=1.5\), \(\ell_{\mathrm{cut}}=8\) yields \(\mathcal{M}_{\mathrm{cpl}}\lesssim10^{-2}\) of total coupling mass
\cite{hansen1988,harrington2001}.

\paragraph{Problem (coupling versus DOS).}
\textbf{Problem.} Channel DOS grows as \(\ell_{\max}^{2}\) but measured diversity saturates. Reconcile.

\textbf{Step 1.} Distinguish listable channels from couplings above \(\varepsilon\sigma_{1}\) \cite{hansen1988,stratton1941}.

\textbf{Step 2.} Apply Theorem~\ref{thm:ch4.coupling-tail-compression} \cite{harrington2001,devaney2004}.

\textbf{Step 3.} Report \(r_{\mathrm{rad}}(\varepsilon)\), not raw DOS \cite{jackson1999}.

\textbf{Physical meaning.} DOS counts labels; coupling tails count watts \cite{hansen1988}.

\paragraph{Problem (MoM channel inflation).}
\textbf{Problem.} MoM uses \(\ell_{\max}=25\); SVD shows \(r_{\mathrm{rad}}(10^{-2})=11\). Explain.

\textbf{Step 1.} Compute \(\mathcal{M}_{\mathrm{cpl}}(\ell_{\mathrm{supp}})\) via Eq.~\eqref{eq:ch4.coupling-tail-mass} \cite{hansen1988}.

\textbf{Step 2.} Identify \(\ell_{\mathrm{supp}}\) from Eq.~\eqref{eq:ch4.ell-max-supported} \cite{harrington2001}.

\textbf{Step 3.} Truncate \(\mathbf{T}\) to \(r_{\mathrm{rad}}\) rows for synthesis \cite{devaney2004,stratton1941}.

\textbf{Physical meaning.} Matrix dimension is coordinate freedom; rank is realizability \cite{harrington2001}.

\paragraph{Problem (anisotropic feed).}
\textbf{Problem.} Offset feed raises \(\mathbf{N}\)-sector coupling at \(\ell=4\) but not \(\mathbf{M}\)-sector. Effect on rank?

\textbf{Step 1.} Evaluate sector-resolved \(\mathcal{M}_{\mathrm{cpl}}^{(\sigma)}\) \cite{harrington2001,hansen1988}.

\textbf{Step 2.} Sum active sectors only for \(r_{\mathrm{rad}}\) \cite{stratton1941}.

\textbf{Step 3.} Report sector ceilings before combined OAM labels \cite{devaney2004}.

\textbf{Physical meaning.} Tail decay is sector-specific; DOS is not \cite{hansen1988}.

\paragraph{Compression optimality link.}
Rank-\(r_{\mathrm{rad}}(\varepsilon)\) truncation is optimal for \(\boldsymbol{\mathcal{K}}_{\mathrm{geo}}\) in operator norm
\cite{harrington2001,devaney2004}. Retaining \(\ell_{\max}\gg\ell_{\mathrm{supp}}\) adds null directions that MoM may misinterpret as
coupled channels \cite{stratton1941,hansen1988}.

\paragraph{DOS inflation diagnostic.}
Define the \emph{listing inflation factor}
\begin{equation}
\label{eq:ch4.dos-inflation}
\mathcal{I}_{\mathrm{DOS}}
=
\frac{(2\ell_{\max}+1)^{2}}{d_{\mathrm{eff}}(\varepsilon)}
\approx
\frac{\text{channel labels}}{\text{export directions}}.
\end{equation}
When \(\mathcal{I}_{\mathrm{DOS}}\gg1\), MoM practitioners are fitting noise in padding rows \cite{harrington2001,devaney2004,hansen1988}.
Theorem~\ref{thm:ch4.coupling-tail-compression} explains why \(\mathcal{I}_{\mathrm{DOS}}\) grows with \(\ell_{\max}\) at fixed \(k_{0}D_{s}\)
\cite{stratton1941,jackson1999}.

\begin{corollary}[Safe listing width]
\label{cor:ch4.safe-listing-width}
A conservative listing satisfies \(\ell_{\max}\le\ell_{\mathrm{supp}}(\varepsilon,k_{0}D_{s})+\mathcal{O}(1)\) in the aligned gauge
\cite{hansen1988,harrington2001}. Wider listings do not enlarge \(\mathcal{S}_{\mathrm{real}}(\mathcal{C})\) \cite{devaney2004,stratton1941}.
\end{corollary}

\paragraph{Worked illustration (inflation factor).}
\(\ell_{\max}=20\), \(k_{0}D_{s}=1.2\), \(d_{\mathrm{eff}}(10^{-2})=9\) gives \(\mathcal{I}_{\mathrm{DOS}}\approx(41)^{2}/9\approx190\)
\cite{hansen1988,harrington2001,devaney2004}. Ninety-nine percent of matrix rows are synthesis-invisible at tolerance \cite{stratton1941}.

\paragraph{Problem (safe \(\ell_{\max}\)).}
\textbf{Problem.} Choose \(\ell_{\max}\) for MoM without padding inflation.

\textbf{Step 1.} Compute \(\ell_{\mathrm{supp}}(10^{-2},k_{0}D_{s})\) \cite{hansen1988,harrington2001}.

\textbf{Step 2.} Set \(\ell_{\max}=\ell_{\mathrm{supp}}+\Delta\) with \(\Delta\le2\) \cite{devaney2004}.

\textbf{Step 3.} Verify \(\mathcal{I}_{\mathrm{DOS}}\lesssim10\) \cite{stratton1941,jackson1999}.

\textbf{Physical meaning.} Listing width should track supported order, not numerical habit \cite{hansen1988}.

\paragraph{Validity.}
Eq.~\eqref{eq:ch4.coupling-tail-envelope} is narrowband and assumes isotropic exterior \cite{stratton1941}. Anisotropic feeds modify
\(\mathcal{C}_{\mathrm{tail}}\) without removing exponential decay \cite{harrington2001}. Loss broadens resonances but does not convert
sub-threshold \(\ell\) sectors into export directions on passive supports \cite{devaney2004,stratton1941}.

